\documentclass[11pt]{article}
\usepackage[margin=1in]{geometry}
\usepackage{amsmath,amssymb,amsthm}
\usepackage{booktabs}
\usepackage{listings}
\usepackage{xcolor}
\usepackage[numbers]{natbib}
\usepackage{hyperref}

\lstdefinelanguage{lean}{
  morekeywords={def,theorem,lemma,example,namespace,end,import,where,instance,fun,by,decide,Prop,Type,open,intro,exact,rw,refine,have,simpa,apply,with,induction,using,by_contra,push_neg,obtain,cases,unfold,omega},
  sensitive=true, morecomment=[l]{--}, morecomment=[s]{/-}{-/}, morestring=[b]"
}
\lstset{language=lean, basicstyle=\ttfamily\small, breaklines=true, columns=fullflexible,
  keywordstyle=\color{blue!60!black}\bfseries, commentstyle=\color{gray},
  literate={∀}{{$\forall$}}1 {→}{{$\to$}}1 {↔}{{$\leftrightarrow$}}1 {∧}{{$\land$}}1
    {∨}{{$\lor$}}1 {¬}{{$\neg$}}1 {∈}{{$\in$}}1 {≤}{{$\le$}}1 {∏}{{$\prod$}}1
    {ℝ}{{$\mathbb{R}$}}1 {ℕ}{{$\mathbb{N}$}}1 {ι}{{$\iota$}}1 {·}{{$\cdot$}}1 {−}{{$-$}}1
    {⊤}{{$\top$}}1 {⊥}{{$\bot$}}1 {⟨}{{$\langle$}}1 {⟩}{{$\rangle$}}1 {←}{{$\leftarrow$}}1
    {⦃}{{$\{$}}1 {⦄}{{$\}$}}1 {⊢}{{$\vdash$}}1 {∃}{{$\exists$}}1}

\newtheorem{theorem}{Theorem}
\newtheorem{lemma}[theorem]{Lemma}
\newtheorem{proposal}{Proposed statement}

\title{One spare corridor is the least price of surviving\\ any single corridor loss:\\ a note for corridor planners with a Lean proof}
\author{Justin Hart\\Viridis LLC}
\date{25 November 2026}

\begin{document}
\maketitle

\begin{abstract}
Hierarchical Dendritic Forest Management (HDFM) designs tree-shaped (dendritic) corridor networks
between habitat patches. A previous note in this series treated \texttt{P6.MinimalNetwork}: a connected
network on $n$ patches needs at least $n-1$ corridors, it has exactly $n-1$ only when it is a tree, and
then every corridor is a bridge whose loss disconnects the network. This note treats the companion sealed
statement \texttt{P6.SingleFailureRobust}. If a network on $n\ge 2$ patches is connected and stays
connected after the loss of any one corridor, it has at least $n$ corridors. The triangle and the
4-cycle attain the bound. So one corridor beyond a tree is the least price of surviving any single loss.
That one corridor must also be placed well: we observe, with a short informal argument that we did not
formalise, that a network with exactly $n$ corridors survives every single loss only when it is a single
cycle through all patches. A corridor added to a star, for example, leaves most corridors as bridges. We give a Lean~4
proof of the statement against the pinned Mathlib, by contradiction from Mathlib's tree
characterisation, and we prove the two attainment witnesses in the same file. These are classical facts
of finite graph theory, not new theorems. The graph carries no geometry, width, quality or cost, and the
result says nothing about which corridors a real landscape should add.
\end{abstract}

\section{Introduction}
Graph models of landscape connectivity represent habitat patches as vertices and functional
links between them as edges~\citep{urbankeitt2001}. Urban and Keitt argued that the minimum spanning
tree is a useful guide to which patches and links matter most to connectivity~\citep{urbankeitt2001}.
HDFM goes further and \emph{designs} dendritic corridor networks. Such a design is as cheap as a
connected network can be. The previous session (P6-T002) proved \texttt{P6.MinimalNetwork}, which shows
the cost of that economy: in a tree every corridor is a bridge.

\paragraph{Question.} What is the least number of corridors that lets a network on $n$ patches survive
the loss of any one corridor, and is that number attained?

\paragraph{Answer.} At least $n$ corridors are needed when $n\ge 2$, one more than a tree has. The triangle
($n=m=3$) and the 4-cycle ($n=m=4$) show the bound is attained. Section~\ref{sec:placement} explains why the
placement of the extra corridor matters as much as its existence.

Justin Hart sealed \texttt{P6.SingleFailureRobust} by decision D-025. This session (P6-T003) proves it.

\section{Background and related work}
A graph is \emph{2-edge-connected} if it is connected and has no bridge; that is exactly the hypothesis
\texttt{SurvivesSingleFailure} on graphs with at least two vertices. Every edge of a 2-edge-connected graph
lies on a cycle, and a graph is acyclic exactly when every edge is a bridge. A connected graph on $n$
vertices has at least $n-1$ edges, with equality exactly for trees~\citep{diestel2017}. The bound $m\ge n$
follows at once. We use Lean~4~\citep{lean4} and Mathlib~\citep{mathlib2020}, which states the needed facts
with \texttt{Nat.card}: \texttt{SimpleGraph.Connected.card\_vert\_le\_card\_edgeSet\_add\_one},
\texttt{SimpleGraph.isTree\_iff\_connected\_and\_card},
\texttt{SimpleGraph.isAcyclic\_iff\_forall\_edge\_isBridge} and \texttt{SimpleGraph.isBridge\_iff}.
This note contributes no new mathematics. It contributes a sealed statement whose hypotheses are listed
in full, a proof that a machine can check, and a plain account of what that proof does and does not cover.

\section{The model and the statements}
Patches form a finite type \texttt{V : Type} with decidable equality. Corridors are the edges of a
\texttt{G : SimpleGraph V} with decidable adjacency. We write $n=\texttt{patchCount V}=\texttt{Fintype.card V}$ and
$m=\texttt{corridorCount G}=\texttt{G.edgeFinset.card}$. Removing corridor $e$ is \texttt{G.deleteEdges \{e\}}.
A simple graph has no loops and at most one corridor between two patches. The model has no geometry, width,
quality or cost; a corridor either exists or does not.

\begin{theorem}[\texttt{P6.MinimalNetwork}, proved in P6-T002]\label{thm:min}
For every finite $V$ and simple graph $G$ on $V$: (1) if $G$ is connected then $n\le m+1$; (2) $G$ is
connected with $m+1=n$ if and only if $G$ is a tree; (3) if $G$ is a tree then $G-e$ is disconnected for
every edge $e$.
\end{theorem}

\begin{theorem}[\texttt{P6.SingleFailureRobust}]\label{thm:robust}
For every finite $V$ with $n\ge 2$ and every simple graph $G$ on $V$: if $G$ is connected and $G-e$ is
connected for every edge $e$ of $G$, then $n\le m$.
\end{theorem}

The sealed Lean definitions are:
\begin{lstlisting}
def SurvivesSingleFailure {V : Type} (G : SimpleGraph V) : Prop :=
  G.Connected ∧ ∀ e ∈ G.edgeSet, (G.deleteEdges {e}).Connected

def Statement : Prop :=
  ∀ (V : Type) [Fintype V] [DecidableEq V] (G : SimpleGraph V) [DecidableRel G.Adj],
    2 ≤ patchCount V → SurvivesSingleFailure G → patchCount V ≤ corridorCount G
\end{lstlisting}
The module restates the counting definitions rather than importing \texttt{P6.MinimalNetwork}, so the
proof below reproves the two Mathlib bridges it needs.

\paragraph{The hypothesis $n\ge 2$.} On one patch the empty network is connected and has no corridor to
lose, yet $m=0<1=n$; the module checks this negative control by \texttt{decide}. For $n=2$ no simple graph
satisfies the hypotheses (the only connected one is a single corridor, which is a bridge), so the
statement holds there vacuously. The bound is informative from $n=3$ on.

\section{Proof}
\paragraph{Mathematical proof.} Suppose $G$ satisfies the hypotheses but $m<n$. Since $G$ is connected,
$n\le m+1$ by Theorem~\ref{thm:min}(1), so $m+1=n$ and $G$ is a tree by Theorem~\ref{thm:min}(2). Since
$n\ge 2$ there are distinct patches $u\neq w$, and the first step of a walk from $u$ to $w$ in $G$ is a
corridor $e$. By Theorem~\ref{thm:min}(3), $G-e$ is disconnected, contradicting the hypothesis.

\paragraph{Lean proof.} The Lean proof follows the same route. It translates \texttt{Fintype.card} and
\texttt{edgeFinset.card} into \texttt{Nat.card}, so it works for arbitrary instances. It obtains $m+1=n$ by
\texttt{omega}, gets the tree from \texttt{isTree\_iff\_connected\_and\_card}, finds an edge by case analysis
on a walk between two distinct vertices, and reuses the bridge lemma from the P6-T002 proof. The same file
proves both attainment witnesses by \texttt{decide}, after unfolding \texttt{SurvivesSingleFailure}. The
witnesses are separate theorems: the sealed \texttt{Statement} is only the bound. The complete candidate
file follows.

\begin{lstlisting}
import Mathlib
import ViridisStatements.P6.SingleFailureRobust

namespace Candidate

open SimpleGraph ViridisStatements.P6.SingleFailureRobust

/-- The number of corridors as a `Nat.card` of the edge set. -/
lemma corridorCount_eq {V : Type} [Fintype V] [DecidableEq V] (G : SimpleGraph V)
    [DecidableRel G.Adj] : corridorCount G = Nat.card G.edgeSet := by
  rw [corridorCount, Nat.card_eq_fintype_card, ← edgeFinset_card]

lemma patchCount_eq (V : Type) [Fintype V] : patchCount V = Nat.card V := by
  rw [patchCount, Nat.card_eq_fintype_card]

/-- In a tree, every corridor is a bridge, so deleting it disconnects the network. -/
lemma tree_delete_not_connected {V : Type} (G : SimpleGraph V) (hT : G.IsTree)
    (e : Sym2 V) (he : e ∈ G.edgeSet) : ¬ (G.deleteEdges {e}).Connected := by
  have hb : G.IsBridge e := isAcyclic_iff_forall_edge_isBridge.mp hT.IsAcyclic he
  induction e using Sym2.ind with
  | _ u v =>
    rw [isBridge_iff] at hb
    intro hc
    apply hb.2
    have := hc.preconnected u v
    simpa [deleteEdges] using this

/-- A connected graph on at least two vertices has an edge. -/
lemma exists_edge {V : Type} [Fintype V] (G : SimpleGraph V) (hc : G.Connected)
    (h2 : 2 ≤ Fintype.card V) : ∃ e, e ∈ G.edgeSet := by
  obtain ⟨u, v, huv⟩ := Fintype.exists_pair_of_one_lt_card (by omega : 1 < Fintype.card V)
  obtain ⟨p⟩ := hc.preconnected u v
  cases p with
  | nil => exact absurd rfl huv
  | cons h _ => exact ⟨_, (G.mem_edgeSet).mpr h⟩

theorem main : ViridisStatements.P6.SingleFailureRobust.Statement := by
  intro V _ _ G _ h2 hS
  obtain ⟨hc, hdel⟩ := hS
  by_contra hlt
  push_neg at hlt
  have hle : patchCount V ≤ corridorCount G + 1 := by
    rw [patchCount_eq, corridorCount_eq]
    exact hc.card_vert_le_card_edgeSet_add_one
  have heq : Nat.card G.edgeSet + 1 = Nat.card V := by
    have h1 := corridorCount_eq G
    have h3 := patchCount_eq V
    omega
  have hT : G.IsTree := isTree_iff_connected_and_card.mpr ⟨hc, heq⟩
  obtain ⟨e, he⟩ := exists_edge G hc h2
  exact tree_delete_not_connected G hT e he (hdel e he)

/-- Attainment witness 1 (restated from the sealed module): the triangle. -/
theorem triangle_attains : 2 ≤ patchCount (Fin 3) ∧
    SurvivesSingleFailure (⊤ : SimpleGraph (Fin 3)) ∧
    corridorCount (⊤ : SimpleGraph (Fin 3)) = patchCount (Fin 3) := by
  unfold SurvivesSingleFailure
  decide

/-- Attainment witness 2 (restated from the sealed module): the 4-cycle. -/
theorem cycle4_attains : 2 ≤ patchCount (Fin 4) ∧
    SurvivesSingleFailure cycle4 ∧
    corridorCount cycle4 = patchCount (Fin 4) := by
  unfold SurvivesSingleFailure
  decide

end Candidate
\end{lstlisting}

\section{Results}
The candidate \texttt{proofs/P6.SingleFailureRobust.lean} compiles under Lean \texttt{v4.28.0} with Mathlib
commit \texttt{8f9d9cff6bd728b17a24e163c9402775d9e6a365}. It passed the local check script, which runs the
denylist, compile, statement-match, axiom and kernel-replay stages. Its only axioms are \texttt{propext},
\texttt{Classical.choice} and \texttt{Quot.sound}. The independent verifier decides whether the claim is
labelled proven; its certification page precedes this paper. The first compile attempt failed on two small
points, recorded in the lab notebook: a walk step yields an adjacency rather than edge-set membership, and
a rewrite with the patch-count bridge also matched the edge count (both are \texttt{Fintype.card}).

\begin{table}[h]
\centering
\begin{tabular}{lccl}
\toprule
Network & $n$ & $m$ & Survives any single loss?\\
\midrule
Any tree (Theorem~\ref{thm:min}) & $n$ & $n-1$ & no (every corridor a bridge)\\
Triangle $K_3$ & 3 & 3 & yes; attains $m=n$ (Lean)\\
4-cycle $C_4$ & 4 & 4 & yes; attains $m=n$ (Lean)\\
Star on 4 patches plus one leaf--leaf corridor & 4 & 4 & no (one corridor is a bridge)\\
Path on 4 patches plus end--end corridor ($=C_4$) & 4 & 4 & yes\\
\bottomrule
\end{tabular}
\caption{Corridor counts and single-failure survival. The first three rows follow from the Lean
development. The fourth and fifth rows are by inspection and illustrate Section~\ref{sec:placement}.}
\label{tab:examples}
\end{table}

\section{One spare corridor, and where it goes}\label{sec:placement}
Read together, the two statements bracket the price of robustness. \texttt{P6.MinimalNetwork} says the
cheapest connected plan is a tree and that none of its corridors is spare. \texttt{P6.SingleFailureRobust}
says that surviving any single corridor loss costs at least one corridor more. The witnesses show that one
more can be enough.

One more is not always enough, though. The following observation is classical. We give it with an informal
argument and have \emph{not} formalised it. Suppose $G$ is connected with $m=n\ge 3$. Then $G$ contains
exactly one cycle, since a spanning tree has $n-1$ edges and the single extra edge closes exactly one
cycle. An edge lies on no cycle exactly when it is a bridge. So $G$ survives every single loss exactly when
every edge lies on that one cycle, that is, when $G$ \emph{is} a cycle through all $n$ patches. A
planner who starts from a dendritic network and buys one extra corridor therefore gains full single-failure
robustness only if the tree is a path and the new corridor joins its two ends. Any other placement protects
only the corridors on the new cycle and leaves the rest as bridges. In a star with a corridor added between two
leaves (Table~\ref{tab:examples}), the corridor to the third leaf is still a bridge. The bound is about how
many corridors; whether a plan is robust also depends on where they go.

We make no claim about which corridors a real landscape should add. Real corridor networks are not chosen
from abstract graphs: cost, distance, species and the failure mode all matter, and a single cycle through all
patches may be impractical or useless for a given species.

\section{Limitations}
\begin{itemize}
\item These are classical facts about finite graphs, not new theorems.
\item The graph carries no geometry, width, quality or cost. The result is about the \emph{number} of
corridors, not their length or price. Patch (vertex) losses, simultaneous losses of two or more corridors, and
minimum-cost robust networks are not covered.
\item The placement observation of Section~\ref{sec:placement} is argued informally and is not part of any sealed
statement.
\item ``Connected'' means connected in the abstract graph. A proof certifies the mathematics of the model
only. It does not certify that a mapped corridor carries movement, that corridor losses happen one at a time, or
any ecological outcome.
\item Multigraphs are outside the model. Two parallel corridors between the same pair of patches would survive
a single loss with $m=2=n$, which a simple graph cannot do; the theorem says nothing about such cases.
\item It does not discharge the FaceCount assumption of \emph{Trees Cannot Cut} (project P1).
\end{itemize}

\section{Reproducibility}
The candidate proof is \texttt{proofs/P6.SingleFailureRobust.lean}. It imports only \texttt{Mathlib} and the
sealed module \texttt{ViridisStatements.P6.SingleFailureRobust} (sha256 \texttt{f7c92557\ldots 2c81}) and uses
none of the forbidden constructs. To check it, run
\texttt{viridis-lean-check proofs/P6.SingleFailureRobust.lean P6.SingleFailureRobust} under toolchain
\texttt{leanprover/lean4:v4.28.0} and Mathlib commit \texttt{8f9d9cff6bd728b17a24e163c9402775d9e6a365}.
The P6-T002 proof of \texttt{P6.MinimalNetwork} is in that run's package. No data or numerical code are
involved in this note.

\section*{Changelog}
\begin{itemize}
\item 2026-11-23 (P6-T001): first version. Drafted three statement proposals
(\texttt{P6.MinimalNetwork}, \texttt{P6.SingleFailureRobust}, \texttt{P6.SerialParallel}).
\item 2026-11-24 (P6-T002): \texttt{P6.MinimalNetwork} sealed (D-024); paper rewritten as a note on it,
with mathematical and Lean proofs and the planning trade-off.
\item 2026-11-25 (P6-T003): \texttt{P6.SingleFailureRobust} sealed (D-025). Paper refocused on it: statement,
the role of $n\ge 2$, mathematical and Lean proofs, attainment witnesses, an example table, and an informal
section on placement of the extra corridor. The \texttt{P6.MinimalNetwork} proof listing was replaced by its
statement; it remains in the P6-T002 package.
\end{itemize}

\bibliographystyle{plainnat}
\bibliography{refs}
\end{document}
